\subsection{Rationale for ABE-Backed Service Authorization}
\label{sec:abe-rationale}

NDNSF uses ABE to protect discovery descriptors and reuses attribute-based
service permissions in its invocation checks. The first choice addresses
confidentiality before the participating Providers are known; the second
keeps service entitlements independent of identity-signing credentials.
Neither implies that ABE is always preferable to role-based authorization
with suitable encryption. Withdrawal and rekeying are lifecycle costs to evaluate.
Execution authorization determines whether a User may invoke a service;
Provider authorization determines whether an entity may offer that service.
Content confidentiality limits who can read protected content. The Controller
administers these service permissions; the User and Provider check them during
invocation. Signature validation under trust rules and invocation-state checks
remain necessary. Permission to offer a service is distinct from selection for
one particular task.

\paragraph{Reason 1: confidential discovery before Provider selection.}
Before Provider selection, the User publishes a discovery descriptor containing
the location, deadline, and resource requirements needed for a Provider's
positive-ACK decision. For example, nearby UAVs need an observation area and
arrival deadline to decide whether to participate, but do not yet need the full
application input. The User encrypts this descriptor using ABE and the requested
service's attribute, without enumerating Provider identities. Confidentiality is against parties
lacking the required decryption capability: authorized but unselected Providers
can still read the descriptor. Content encryption does not hide outer Data names
or traffic patterns. In the request-scoped confidentiality path, the full
application input is encrypted with a request-scoped key; Selection protects
that key for the selected Provider. The requesting User retains access. A Trust Schema alone does
not encrypt content; a role-authorized design can also use ABE or service-group
encryption to provide discovery confidentiality.

Encrypted Request Data can be retrieved or cached without granting decryption
access. This separation of delivery from disclosure also holds for other
content-encryption schemes; it is not unique to ABE.

\paragraph{Reason 2: reusing attribute-based permissions for service invocation.}
Confidential discovery alone does not justify ABE-backed execution authorization:
ABE encryption could instead be combined with role-based invocation checks.
NDNSF chooses to administer both Provider and User service permissions through
the existing attribute-authority infrastructure, separately from identity-signing
credentials.

Provider permissions determine access to the Requester Challenge; User permissions
determine access to the Provider Challenge. NDNSF checks the returned challenge
values together with signatures, current permissions, and request state
(Section~\ref{sec:authorization-withdrawal}). Decryption is therefore one part of
the invocation check, not permission to execute any request.

KP-ABE associates attributes with ciphertext and an access policy with each
decryption key (DKEY)~\cite{dulal2023nacabe}. NDNSF uses
\texttt{/SERVICE/<service>} for Provider access and
\texttt{/PERMISSION/<service>} for User access. Within one authority and parameter
generation, an entity's service permissions can be combined in one individually
issued DKEY. For example, one User's policy can allow detection OR mapping,
while another allows detection only. Their identity-signing credentials need
not change to represent these different service sets; each service retains its
own ciphertext attribute.

DNMP instead uses role signing keys and a Trust Schema to constrain command
targets and actions~\cite{dnmp}. In the DNMP-inspired alternative
(Section~\ref{sec:dnmp-alternative}), retaining ABE for discovery would require
managing both decryption permissions and role-based invocation permissions.
NDNSF uses the attribute-authority infrastructure for both service-level checks,
without adding role signing credentials solely to represent service membership.
Trust rules for validating Data signatures remain necessary.

This is a management-reuse rationale, not a demonstrated cost advantage.
Provider and User permissions remain distinct; discovery encryption does not
already supply User permission DKEYs or challenge processing. DKEY issuance,
policy updates, and revocation still incur costs. DNMP roles and signing
credentials can also cover several permissions. For policies requiring
command-argument restrictions, its schema-based checks may be a better fit than
NDNSF's service-level checks. The evaluation must compare both designs at the
same policy granularity, including credential updates and invocation overhead.

\subsection{Authorization Mechanisms, Lifecycle, and Costs}
\label{sec:authorization-withdrawal}

\paragraph{Connecting decryption capability to invocation authorization.}
In the ABE-backed invocation path, the User includes a Requester Challenge
(\texttt{UserToken}) in protected Request content. A Provider that recovers it
returns it in ACK Data; the User checks it against the active request.
The Provider places its own Challenge (\texttt{ProviderToken}) in protected
ACK content. The Provider checks the corresponding challenge response in
Selection before starting the selected task.

These checks demonstrate access to protected challenge values, not the creation
of new permissions or a formal proof of exclusive key possession. The User and
Provider also validate message signatures and certificate chains under trust
rules, check current service permissions, and match the request and its state.
One-time state prevents a repeated message from authorizing another execution
only while that state is preserved or safely invalidated on restart. Authorized
key or challenge sharing is outside the claimed threat protection. This reuses
attribute-based permission management rather
than a separate set of role-based invocation rules. It is not cost-free: confidential
discovery already needs Provider DKEYs, whereas User permission DKEYs and
challenge processing are additional authorization costs. ABE confidentiality
combined with DNMP-inspired role authorization remains a valid alternative.

\paragraph{Runtime permission withdrawal and evaluation.}

The current runtime implements permission withdrawal through Controller-signed
authorization status and checks at invocation processing points. A participant
must retrieve and validate updated status before enforcing the newly withdrawn
authority. ABE parameter rotation and refreshed DKEYs protect subsequent
exchanges encrypted under the new generation. The present rotation is
Controller-wide, so other authorized members may also need refreshed DKEYs;
it is not a zero-fan-out revocation scheme.

DNMP specifies validity bounds and replay-state handling~\cite{dnmp}; these
do not by themselves provide immediate withdrawal of an unexpired role key.
For the comparison, a DNMP-inspired design can revoke
the relevant credential and require execution nodes to check signed revocation
status or reject expired certificates. This invalidates authority at the
validator; it does not physically retrieve distributed certificates.
Both approaches require update propagation, and neither retracts previously
disclosed keys or plaintext.

Focused integration tests exercise NDNSF's runtime revocation checks, including
identity-, certificate-, and service-scoped withdrawal. These tests establish
behavior at the exercised runtime boundaries, not network-wide instantaneous
revocation or lower cost than certificate-based alternatives. The remaining
evaluation should measure status-propagation delay, rejection delay, rekeying
traffic, and effects on unaffected members. Runtime status validation and key
rotation therefore extend NDNSF to permission withdrawal, with explicit
propagation and rekeying costs.
